\begin{lemma}[Terminal-prefix repackaging of a good hierarchy]
Let $G$ be a finite loopless multigraph with distinguished edge $e_0$
and palette $[k]$.  Let $\varphi_{\rm in},\varphi_{\rm gen},\varphi$
be partial edge colorings, and let $\mathcal E$ be a positive-rung
\noderef{ExtendedTashkinovTree} with input coloring
$\varphi_{\rm in}$, generating coloring $\varphi_{\rm gen}$, exact
terminal tree $T$, and \noderef{MaximumProperty}.  Let $\mathcal B$ be
a \noderef{HierarchyBase} under $\varphi$, let $H$ be a
\noderef{TashkinovHierarchy} with divider count $q$, and let $\Gamma$
be a \noderef{GoodHierarchy} on $H$.  Put
\[
                    A=\noderef{HierarchyStarredStage}(H,q).
\]
Suppose that $B$ is an exact ordered tree-sequence prefix satisfying
$A\prec B\preccurlyeq T$, where both comparisons retain every displayed
vertex and individual edge identity as in
\noderef{TreeSequenceInitialSegment}.

Then there exist an ETT $\mathcal E_B$ on the actual tree $B$, a base
$\mathcal B_B$, a hierarchy $H_B$, and a good hierarchy $\Gamma_B$
such that all of the following hold.
\begin{enumerate}
\item $\mathcal E_B$ has the same rung number and the same complete
certified series as $\mathcal E$, and has maximum property.
\item The core, closure, and numerical support position of
$\mathcal B_B$ are those of $\mathcal B$.
\item $H_B$ has divider count $q$.  Its stages at indices
$0,\ldots,q$, its starred stages, and all numerical divider positions
are exactly the corresponding data of $H$; only its terminal stage is
changed from $T$ to $B$.
\item For every $0\leq j\leq q$, the computed stored-color and
stage-missing-color sets agree:
\[
 D^{H_B}_{n,j}=D^H_{n,j},\qquad
 \noderef{HierarchyStageMissingColors}(H_B,j)
  =\noderef{HierarchyStageMissingColors}(H,j).
\]
For every color $\eta$, the Gamma data agree and the new segment-color
set is contained in the old one:
\[
 \Gamma^j_{B,\eta}=\Gamma^j_\eta,\qquad
 \noderef{HierarchySegmentColors}(H_B,j,\eta)
 \subseteq \noderef{HierarchySegmentColors}(H,j,\eta).
\]
\end{enumerate}
All equalities here concern the complete ordered edge-valued sequences,
not merely their vertex sets.
\end{lemma}
\begin{proof}
Write the nonterminal stages of $H$ as
$T_0,T_1,\ldots,T_q$ and its terminal stage as $T_{q+1}=T$.
The nestedness and terminal clauses of
\noderef{TashkinovHierarchy} show, by composition, that every $T_j$ is
an exact initial segment of $T$.  If $q>0$, the definition
\noderef{HierarchyStarredStage} gives $A=T_q$.  If $q=0$, it gives
$A=\mathcal B.\mathit{core}$.  In the latter case, the non-PE clause of
\noderef{HierarchyBase} identifies this core with $T_0$, while its PE
clause supplies an exact \noderef{OrderedTreeUnion} in which $T_0$ is
an ordered initial segment of the core.  Thus in both cases
the last ladder tree $T_0$ is an initial segment of $A$ and hence of
$B$.  It is strictly shorter than $B$: this follows through
$T_0\preccurlyeq A\prec B$ when $q=0$, and from the strict hierarchy
stages followed by $A=T_q\prec B$ when $q>0$.

Construct $\mathcal E_B$ by retaining the rung number and the entire
certified series of $\mathcal E$ and replacing only its displayed tree
$T$ by $B$.  The preceding paragraph supplies its strict lower
containment.  Composing $B\preccurlyeq T$ with the old upper containment
of $T$ in the final certified closure supplies its upper containment.
The generating-coloring equation is unchanged.  These are every field
of \noderef{ExtendedTashkinovTree}.  Since
\noderef{MaximumProperty} depends only on the retained rung number and
certified series, its old witness proves maximum property for
$\mathcal E_B$.

The series, its terminal ladder tree, and its terminal kind did not
change.  Copy the core, closure, support position, support positivity,
terminal missing-color equality, PE clause, and non-PE clause of
$\mathcal B$.  Each copied term has exactly the required type, so these
data construct the asserted \noderef{HierarchyBase} $\mathcal B_B$.

Define $H_B$ with divider count $q$, stages $T_0,\ldots,T_q,B$, and the
old divider positions.  The base stage, all retained nestedness and
strictness clauses, and divider positivity and increase copy from $H$.
If $q>0$, the first ordinary stage is retained, so core-inside-first
also copies.  If $q=0$, the first stage of $H_B$ is its terminal stage
$B$, and core-inside-first is the assumed relation $A\preccurlyeq B$
because $A$ is the copied core.  The only new nesting and
strictness clause is $T_q\prec B$ when $q>0$, which follows from
$T_q=A\prec B$; when $q=0$ it is $T_0\prec B$, proved above.  For a
divider $i<q$, the old divider-stage identity writes $T_{i+1}$ as the
prefix of $T$ at the old numerical divider position.  Both
$T_{i+1}$ and $B$ are exact prefixes of $T$, and
$T_{i+1}\preccurlyeq T_q=A\prec B$, so entrywise comparison in
\noderef{TreeSequenceInitialSegment} identifies $T_{i+1}$ with the
prefix of $B$ at that same position.  Hence every divider-stage identity
also holds.  This checks every field of
\noderef{TashkinovHierarchy} and proves item 3.

We isolate the totalized first-missing argument needed for goodness.
Let $P\preccurlyeq Q$ be exact ordered tree sequences, let $\eta$ be a
color, and let $s<|P|$ satisfy
\[
 s<\noderef{HierarchyFirstMissingPosition}(\varphi,P,\eta).
                                                               \tag{1}
\]
Then $s$ is smaller than the first missing position in $Q$.  There are
two cases in the definition.  If $Q$ has no vertex missing $\eta$, its
totalized position is its last index $|Q|$, and $s<|P|\leq|Q|$.  If it
has a missing candidate and its least candidate $t$ were at most $s$,
then $t<|P|$.  Exact prefix equality would make the vertex of $P$ at
$t$ miss $\eta$.  If $P$ has a missing candidate, its least candidate
would be at most $t\leq s$, contradicting (1); if $P$ has none, the
candidate at $t$ is already a contradiction.  This proves
\[
 s<\noderef{HierarchyFirstMissingPosition}(\varphi,Q,\eta).   \tag{2}
\]
The argument explicitly includes both no-candidate fallback branches.

Apply (1)--(2) to $P=B$ and $Q=T$.  Fix $j\leq q$, a color $\eta$,
and a color $c$ in the $(j,\eta)$ segment-color set of $H_B$.  By
\noderef{HierarchySegmentColors}, there is an exact edge position $s$
above the common starred stage, with
$s<|H_B.\mathit{stage}(j+1)|\leq|B|$, such that the displayed edge has
color $c$ and $s$ is below the first $\eta$-missing position in $B$.
The preceding observation puts $s$ below the first $\eta$-missing
position in $T$.  When $j<q$, the next stage is retained literally.
When $j=q$, the new next stage is $B$ and the old one is $T$; the exact
prefix relation $B\preccurlyeq T$ identifies their edge at $s$ and gives
$s<|T|$.  Thus in both cases the same position, with the same exact edge
and $\varphi$-color, witnesses that $c$ lies in the old $(j,\eta)$
segment-color set.  Therefore
\[
 \noderef{HierarchySegmentColors}(H_B,j,\eta)
 \subseteq\noderef{HierarchySegmentColors}(H,j,\eta).         \tag{3}
\]
Only this direction is proved or needed; in particular, no equality of
the two full segment-color functions is asserted.

Define $\Gamma_B^j(\eta)=\Gamma^j(\eta)$.  The corresponding raw
certified stored set and starred stage are unchanged, so unfolding
\noderef{HierarchyStoredColorsAt} gives
$D^{H_B}_{n,j}=D^H_{n,j}$.  At $j=0$ the stage-missing set is the
missing-color set of the unchanged terminal ladder tree; at $j>0$ it is
the missing-color set of the unchanged ordinary stage.  Thus
\noderef{HierarchyStageMissingColors} gives the second equality in
item 4.

It remains to verify all six fields of
\noderef{GoodHierarchy}.  The two-color cardinality and pairwise
disjointness fields are the old fields after the stored-set and Gamma
equalities.  For avoidance, an old Gamma color lies in the common
stage-missing set and outside the old segment-color set; containment
(3) puts it outside the new segment-color set as required.  In every
unique-refresh clause, both ordinary stages, the preceding starred
stage, their exact vertex-set difference, the stored sets, and all Gamma
pairs are unchanged, so the old unique witness and its uniqueness proof
copy verbatim.  The parallel-base clause uses only the unchanged series,
core, closure, zeroth stored set, zeroth Gamma union, and their
$\varphi$-missing sets, so both its non-minus-closed and cardinality
parts copy.  Finally, each positive-stage minus-closed clause concerns
an unchanged ordinary stage, stored set, and preceding Gamma union, and
is exactly the corresponding old clause.  These six checks construct
$\Gamma_B$ and prove every assertion of the lemma.
\end{proof}
